\documentclass[../cours_tle.tex]{subfiles}
\begin{document}
\setcounter{chapter}{1}
\chapter{Récurrence et sommes}
\label{ch:2}
\begin{objectif}
Spécialité. Prouver une affirmation pour tous les entiers. Prérequis : chapitre 1 et suites de première.
\end{objectif}
\section{Du calcul à la preuve}
Une suite est une fonction de $\N$ (ou d'une partie de $\N$) dans $\R$. Une formule explicite donne $u_n$ directement ; une relation de récurrence exige aussi une ou plusieurs valeurs initiales.
\begin{theoreme}[Principe de récurrence]
Si $P(n_0)$ est vraie et si, pour tout $n\ge n_0$, $P(n)\Rightarrow P(n+1)$, alors $P(n)$ est vraie pour tout $n\ge n_0$.
\end{theoreme}
C'est un fondement des entiers naturels, admis ici. L'hérédité seule ne suffit pas : « tous les entiers sont supérieurs à $100$ » se transmet d'un entier au suivant mais ne s'initialise pas à $0$.
\begin{exemple}[Une suite définie par itération]
$u_0=0$ et $u_{n+1}=\sqrt{2+u_n}$. Montrons $0\le u_n\le2$.
C'est vrai au rang $0$. Sous cette hypothèse, $2\le2+u_n\le4$, donc $0\le u_{n+1}\le2$.
Cela prouve à la fois que chaque terme existe et qu'il est borné.
\end{exemple}
Une récurrence forte suppose $P(n_0),\ldots,P(n)$ pour établir $P(n+1)$. Elle se ramène à la récurrence ordinaire appliquée à leur conjonction.
Pour une relation d'ordre $2$, deux valeurs initiales sont généralement nécessaires.
\section{Sommes arithmétiques, géométriques et télescopiques}
\begin{propriete}
Pour $n\ge1$, $\sum_{k=1}^n k=n(n+1)/2$. Pour $q\ne1$,
$\sum_{k=0}^n q^k=(1-q^{n+1})/(1-q)$ ; pour $q=1$, la somme vaut $n+1$.
\end{propriete}
\begin{demonstration}
Pour la première somme, initialiser à $1$ ; ajouter $n+1$ à $n(n+1)/2$ donne $(n+1)(n+2)/2$.
Pour la seconde, noter $S=1+q+\cdots+q^n$. Soustraire $qS$ à $S$ annule tous les termes intermédiaires et donne $(1-q)S=1-q^{n+1}$.
\end{demonstration}
Le symbole $\sum$ comporte une variable muette : $\sum_{k=1}^n k=\sum_{j=1}^n j$. Il faut compter les termes : de $k=0$ à $n$, il y en a $n+1$.
Une somme télescopique annule des termes successifs :
\[
\sum_{k=1}^n\frac1{k(k+1)}=\sum_{k=1}^n\left(\frac1k-\frac1{k+1}\right)=1-\frac1{n+1}.
\]
\section{L'inégalité de Bernoulli}
\begin{theoreme}
Pour $x\ge-1$ et $n\in\N$, $(1+x)^n\ge1+nx$.
\end{theoreme}
\begin{demonstration}
Au rang $0$, égalité. Multiplier l'hypothèse par $1+x\ge0$ conserve le sens :
$(1+x)^{n+1}\ge(1+nx)(1+x)=1+(n+1)x+nx^2\ge1+(n+1)x$.
\end{demonstration}
Cette preuve explique la restriction sur $x$ et servira à étudier $q^n$ sans supposer déjà la limite de l'exponentielle.
\section{Une recherche guidée et un invariant}
Pour conjecturer $\sum_{k=1}^n k^2$, calculer les premiers cas peut suggérer une formule ; cela ne la prouve pas. La formule
$n(n+1)(2n+1)/6$ se vérifie à $n=1$ ; sa différence entre les rangs $n+1$ et $n$ vaut $(n+1)^2$ après développement, ce qui établit l'hérédité.

Dans une boucle qui additionne $1,\ldots,n$, l'invariant « après $k$ passages, la somme vaut $k(k+1)/2$ » est une récurrence sur les passages. Initialisation, conservation et terminaison constituent une preuve du programme.
\begin{attention}
Ne pas supposer la conclusion pour tous les rangs lors de l'hérédité. On la suppose à un rang arbitraire fixé, puis on la démontre au rang suivant. Ne pas remplacer la preuve par une capture d'écran des premiers termes.
\end{attention}
\begin{finchapitre}
\section{Utilité et suite du parcours}
Les récurrences structurent les preuves sur les matrices, les graphes et les algorithmes. Les sommes télescopiques deviendront des outils de séries et de calcul intégral.
\begin{retenir}
Initialiser, formuler une hypothèse précise, prouver la transmission, conclure avec le bon rang de départ.
\end{retenir}
\end{finchapitre}
\end{document}
